\section{Intersection coordinates, endpoint jet, and restart}

The mixed-index and critical-height exhaustions have now identified the same
all-orders membership of the maximal history.  The endpoint argument no longer
distinguishes their coordinates.  For every finite spatial order, the common
family contains the adjacent velocity and pressure rows needed to reach
\(T_*\), recover one endpoint jet, and restart the original history.


\begin{definition}[Compatible exhaustive rectangle family]
\label{u:def:compatible-intersection}
The datum-generated rectangles form the family
\begin{equation}
 \left\{\cR_{N,M}[u_0;T_*]:N,M\in\Nzero\right\},
 \qquad
 \left.\cR_{N,M}[u_0;T_*]\right|_{(N',M')}
 =\cR_{N',M'}[u_0;T_*]
 \quad(N'\le N,\ M'\le M).
 \label{u:eq:compatible-intersection-space}
\end{equation}
It is exhaustive in the direct finite-row sense: for every \(r,m\in\Nzero\),
a rectangle with \(N\ge r\) and \(M\ge m\) contains
\begin{equation}
 V_n\in L^2(0,T_*;H^m),
 \qquad Q_n\in L^1(0,T_*;H^m),
 \qquad 0\le n\le r.
 \label{u:eq:compatible-intersection}
\end{equation}
\end{definition}

\Needspace{14\baselineskip}
\begin{theorem}[Compatible exhaustive rectangle family]
\label{u:thm:compatible-inverse-intersection}

Let \(\nu>0\), \(u_0\in\HsSigma(\R^3)\), and let
\((T_*,u,p)\) be the maximal pressure-complete classical history generated
by \((u_0,\nu)\).  If \(T_*<\infty\), then
\begin{equation}
 \left.\cR_{N,M}[u_0;T_*]\right|_{(N',M')}
 =\cR_{N',M'}[u_0;T_*]
 \qquad(N'\le N,\ M'\le M),
 \label{u:eq:ier-compatible-family}
\end{equation}
and every row is the corresponding derivative of the original
pressure-complete history.  Every pair \((V_n,Q_n)\) belongs to a sufficiently
large rectangle \(\cR_{N,M}[u_0;T_*]\) and retains the memberships
\begin{align}
 V_n&\in\bigcap_{m\in\Nzero}L^2(0,T_*;H^m_\sigma(\R^3)),
 \label{u:eq:ier-inverse-velocity-coordinate}\\
 Q_n&\in\bigcap_{m\in\Nzero}L^1(0,T_*;H^m(\R^3)),
 \label{u:eq:ier-inverse-pressure-coordinate}
\end{align}
The family exhausts every finite pair of temporal and spatial orders.  Its
velocity therefore has the direct membership
\begin{equation}
 u\in\bigcap_{r,m\in\Nzero}H^r(0,T_*;H^m(\R^3)),
 \label{u:eq:mixed-sobolev-intersection}
\end{equation}
and, for every \(r,m\in\Nzero\),
\begin{equation}
 \sum_{n=0}^{r}\norm{\partial_t^nu}_{L^2(0,T_*;H^m)}^2
 \le\mathfrak R_{r,m}^*[u_0,\nu,T_*].
 \label{u:eq:mixed-sobolev-finite-bound}
\end{equation}
\end{theorem}

\begin{proof}
The mixed-index exhaustion in Theorem~\ref{u:thm:datum-intersection} and the
critical-height exhaustion in
Corollary~\ref{u:cor:datum-critical-height-completion}, followed by
Proposition~\ref{u:prop:spatial-column-reconstruction}, give the same velocity
memberships in \eqref{u:eq:ier-inverse-velocity-coordinate} and
\eqref{u:eq:mixed-sobolev-intersection}.  The pressure recurrence then gives
\eqref{u:eq:ier-inverse-pressure-coordinate}.
For each finite \(N,M\), the rectangle is selected directly from the mixed
jet.
Theorem~\ref{u:thm:datum-rectangles} and
Proposition~\ref{u:prop:rectangle-pressure-completion} give its finite row
bounds.  Proposition~\ref{u:prop:rectangle-compatibility} gives
\[
 \left.\cR_{N,M}[u_0;T_*]\right|_{(N',M')}
 =\cR_{N',M'}[u_0;T_*]
 \qquad(N'\le N,\ M'\le M).
\]
Hence the rectangles are compatible.  Every \(V_n,Q_n\) appears in a
sufficiently large rectangle, and every
spatial derivative is fixed by its distributional velocity or pressure
row.  The family consequently exhausts all finite mixed velocity and pressure
rows in either description.

Fix finite \(r,m\), put \(I=(0,T_*)\), and choose \(N\ge r\),
\(M\ge m\). The upper rows of that finite record belong to
\(L^2(I;H^{M+1})\hookrightarrow L^2(I;H^m)\) for \(0\le n\le r\).
Proposition~\ref{u:prop:formal-mixed-recurrences} identifies them as the
distributional derivatives \(V_n=\partial_t^nu\). Thus the finite record
contains every term in the derivative-sum norm of \(H^r(I;H^m)\).
Taking \(N=r\), \(M=m\) gives
\[
 \begin{aligned}
 \norm u_{H^r(I;H^m)}^2
 &=\sum_{n=0}^r\norm{V_n}_{L^2(I;H^m)}^2\\
 &\le\sum_{n=0}^r\norm{V_n}_{L^2(I;H^{m+1})}^2
 \le\mathfrak R_{r,m}^*[u_0,\nu,T_*].
 \end{aligned}
\]
This proves \eqref{u:eq:mixed-sobolev-finite-bound} and shows how each
mixed Sobolev norm is recovered from a finite record.  The mixed-index and
completed-height descriptions agree on that record and hence on the full
compatible family.
\end{proof}

\begin{proposition}[Complete spatial row and common endpoint]
\label{u:prop:ier-complete-spatial-endpoint}
Under the hypotheses of
Theorem~\ref{u:thm:compatible-inverse-intersection}, for every
\(m\in\Nzero\),
\begin{equation}
 u\in W^{1,1}(0,T_*;H^m_\sigma(\R^3))
 \hookrightarrow C([0,T_*];H^m_\sigma(\R^3)).
 \label{u:eq:ier-spatial-row-W11}
\end{equation}
The endpoint values supplied at the different Sobolev orders are the images
of one common state
\begin{equation}
 u_*:=u(T_*)\in\bigcap_{m\in\Nzero}H^m_\sigma(\R^3)
 =H^\infty_\sigma(\R^3).
 \label{u:eq:ier-common-spatial-endpoint}
\end{equation}
\end{proposition}

\begin{proof}
Fix \(m\).  The zeroth temporal coordinate of
\eqref{u:eq:mixed-sobolev-intersection} gives
\[
 u\in L^2(0,T_*;H^{m+2}).
\]
The pressure-completed nonlinearity satisfies
\[
 \norm{\Pp\nabla\!\cdot(u\otimes u)}_{H^m}
 \le C_m\norm u_{H^{m+2}}^2
\]
by Lemmas~\ref{u:lem:sobolev-product} and \ref{u:lem:riesz-leray}.  It belongs
to \(L^1(0,T_*;H^m)\).  This spatial row and \(T_*<\infty\) give
\(\Delta u\in L^1(0,T_*;H^m)\), and the projected equation
\[
 u_t=\nu\Delta u-\Pp\nabla\!\cdot(u\otimes u)
\]
places \(u_t\) in \(L^1(0,T_*;H^m)\).  Also
\(u\in L^1(0,T_*;H^m)\), so the Banach-valued fundamental theorem gives
\eqref{u:eq:ier-spatial-row-W11}.

If \(k\ge m\), the \(H^k\)- and \(H^m\)-continuous representatives are
the original preterminal function after the embedding \(H^k\hookrightarrow
H^m\).  Uniqueness of continuous representatives makes their endpoint
values agree under that embedding.  These compatible values define the
single state in \eqref{u:eq:ier-common-spatial-endpoint}.  Continuity of
divergence from \(H^{m+1}\) to \(H^m\) gives \(\nabla\cdot u_*=0\).
\end{proof}

\begin{corollary}[Full-interval pressure-complete mixed jet]
\label{u:cor:ier-whole-terminal-mixed-jet}
Under the hypotheses of Theorem~\ref{u:thm:compatible-inverse-intersection},
for every \(r,m,n\in\Nzero\) and \(\alpha\in\Nzero^3\),
\begin{align}
 u&\in H^r(0,T_*;H^m_\sigma(\R^3)),
 \label{u:eq:ier-whole-terminal-bochner}\\
 J_{n,\alpha}:=\partial_x^\alpha V_n
 &\in L^2(0,T_*;H^{m+1}(\R^3)),
 \label{u:eq:ier-whole-terminal-mixed-upper}\\
 J_{n+1,\alpha}
 &\in L^2(0,T_*;H^{m-1}(\R^3)),
 \label{u:eq:ier-whole-terminal-mixed-lower}\\
 \Pi_{n,\alpha}:=\partial_x^\alpha Q_n
 &\in L^1(0,T_*;H^m(\R^3)).
 \label{u:eq:ier-whole-terminal-mixed-pressure}
\end{align}
The following finite bounds hold:
\begin{align}
 \sum_{j=0}^{r}\norm{V_j}_{L^2(0,T_*;H^m)}^2
 &\le\mathfrak R_{r,m}^*[u_0,\nu,T_*],
 \label{u:eq:ier-finite-bochner-bound}\\
 \norm{J_{n,\alpha}}_{L^2(0,T_*;H^{m+1})}^2
 +\norm{J_{n+1,\alpha}}_{L^2(0,T_*;H^{m-1})}^2
 &\le\mathfrak R_{n,m+|\alpha|}^*[u_0,\nu,T_*],
 \label{u:eq:ier-finite-mixed-bound}\\
 \norm{\Pi_{n,\alpha}}_{L^1(0,T_*;H^m)}
 &\le C_{m+|\alpha|}^{\rm p}(2^{n+1}-1)
 \mathfrak R_{n,m+|\alpha|}^*[u_0,\nu,T_*].
 \label{u:eq:ier-finite-mixed-pressure-bound}
\end{align}
Every mixed row satisfies the recurrences in
Proposition~\ref{u:prop:formal-mixed-recurrences} on \((0,T_*)\).
\end{corollary}

\begin{proof}
Equations \eqref{u:eq:ier-whole-terminal-bochner} and
\eqref{u:eq:ier-finite-bochner-bound} follow from
Theorem~\ref{u:thm:compatible-inverse-intersection}.  The Fourier multiplier
bound
\[
 \norm{\partial_x^\alpha f}_{H^s}
 \le\norm f_{H^{s+|\alpha|}}
\]
and the rectangle \((n,m+|\alpha|)\) give
\eqref{u:eq:ier-whole-terminal-mixed-upper}--%
\eqref{u:eq:ier-finite-mixed-bound}.  Apply this multiplier bound to the
pressure completion in Proposition~\ref{u:prop:rectangle-pressure-completion}
to obtain \eqref{u:eq:ier-whole-terminal-mixed-pressure} and
\eqref{u:eq:ier-finite-mixed-pressure-bound}.  Proposition
\ref{u:prop:formal-mixed-recurrences} supplies the stated identities.
\end{proof}

\Needspace{24\baselineskip}
\begin{theorem}[Compatible full-interval Sobolev traces]
\label{u:thm:ier-compatible-endpoint-traces}
Assume the hypotheses of
Theorem~\ref{u:thm:compatible-inverse-intersection}.
For \(N,m\in\Nzero\), define
\begin{equation}
 \mathfrak E_{N,m}
 :=(1+T_*^{-1})^{1/2}
 \bigl(\mathfrak R_{N,m}^*[u_0,\nu,T_*]\bigr)^{1/2}.
 \label{u:eq:ier-endpoint-majorant}
\end{equation}
\begin{samepage}
If \(0\le n\le N\) and \(\alpha\in\Nzero^3\), then
\begin{align}
 V_n&\in C([0,T_*];H^m_\sigma(\R^3)),
 &\sup_{0\le t\le T_*}\norm{V_n(t)}_{H^m}
 &\le\mathfrak E_{N,m},
 \label{u:eq:ier-row-trace}\\
 J_{n,\alpha}&\in C([0,T_*];H^m(\R^3)),
 &\sup_{0\le t\le T_*}\norm{J_{n,\alpha}(t)}_{H^m}
 &\le\mathfrak E_{N,m+|\alpha|}.
 \label{u:eq:ier-mixed-row-trace}
\end{align}
There are unique values
\begin{align}
 V_n^*&:=\lim_{t\uparrow T_*}V_n(t)
 \in H^\infty(\R^3),
 \label{u:eq:ier-velocity-endpoint}\\
 J_{n,\alpha}^*&:=\lim_{t\uparrow T_*}J_{n,\alpha}(t)
 =\partial_x^\alpha V_n^*
 \in H^\infty(\R^3),
 \label{u:eq:ier-mixed-velocity-endpoint}
\end{align}
\par
\end{samepage}
such that, for \(k\ge m\), the image of the \(H^k\) trace under
\(H^k\hookrightarrow H^m\) is the \(H^m\) trace.  Quantitatively,
\begin{align}
 \norm{V_n^*}_{H^m}
 &\le\mathfrak E_{n,m},
 \label{u:eq:ier-endpoint-velocity-size}\\
 \norm{V_n(t)-V_n^*}_{H^m}
 &\le(T_*-t)\mathfrak E_{n+1,m},
 \qquad0\le t\le T_*,
 \label{u:eq:ier-endpoint-velocity-modulus}\\
 \norm{J_{n,\alpha}^*}_{H^m}
 &\le\mathfrak E_{n,m+|\alpha|},
 \label{u:eq:ier-endpoint-mixed-size}\\
 \norm{J_{n,\alpha}(t)-J_{n,\alpha}^*}_{H^m}
 &\le(T_*-t)\mathfrak E_{n+1,m+|\alpha|}.
 \label{u:eq:ier-endpoint-mixed-modulus}
\end{align}
 In particular, the zeroth trace agrees with the common endpoint in
 Proposition~\ref{u:prop:ier-complete-spatial-endpoint}, and
\begin{equation}
 u_*=V_0^*\in H^\infty_\sigma(\R^3),
 \qquad
 u\in\bigcap_{a,m\in\Nzero}C^a([0,T_*];H^m_\sigma(\R^3)),
 \qquad
 \partial_t^au(T_*)=V_a^*.
 \label{u:eq:ier-complete-left-velocity-jet}
\end{equation}
\end{theorem}

\begin{proof}
For \(0\le n\le N\), Theorem~\ref{u:thm:compatible-inverse-intersection}
provides a finite rectangle containing the adjacent rows
\[
 V_n\in L^2(0,T_*;H^{m+1}),
 \qquad
 \partial_tV_n=V_{n+1}\in L^2(0,T_*;H^{m-1}).
\]
That rectangle gives both finite norms; the critical-height
description recovers the same coordinates.
With \(A=(1-\Delta)^{1/2}\), the adjacent-row identity is
\[
 \frac{\mathrm d}{\mathrm dt}\norm{V_n(t)}_{H^m}^2
 =2\operatorname{Re}\ip{A^{m-1}V_{n+1}(t)}{A^{m+1}V_n(t)}_{L^2}.
\]
Integrating from the initial time and applying Cauchy--Schwarz gives
\begin{equation}
 \sup_{0\le t\le T_*}\norm{V_n(t)}_{H^m}^2
 \le\norm{V_n(0)}_{H^m}^2
 +2\norm{V_{n+1}}_{L^2(I;H^{m-1})}
    \norm{V_n}_{L^2(I;H^{m+1})},
 \qquad I=(0,T_*).
 \label{u:eq:ier-initial-row-trace-bound}
\end{equation}
The initial norm is finite by Lemma~\ref{u:lem:formal-initial-jet}.
Lemma~\ref{u:lem:adjacent-trace} supplies the continuous representative
on \([0,T_*]\), justifies the identity by Fourier cutoff, and gives the
first membership in
\eqref{u:eq:ier-row-trace} and
\begin{align*}
 \sup_{0\le t\le T_*}\norm{V_n(t)}_{H^m}^2
 &\le T_*^{-1}\norm{V_n}_{L^2H^m}^2
 +2\norm{V_n}_{L^2H^{m+1}}\norm{V_{n+1}}_{L^2H^{m-1}}\\
 &\le(1+T_*^{-1})\mathfrak R_{N,m}^*[u_0,\nu,T_*].
\end{align*}
This proves the bound in \eqref{u:eq:ier-row-trace}.  Apply this lemma to
\(J_{n,\alpha}\), using \eqref{u:eq:ier-finite-mixed-bound}, to prove
\eqref{u:eq:ier-mixed-row-trace}.

For \(k\ge m\), the \(H^k\)-continuous and \(H^m\)-continuous
representatives equal the original preterminal distribution.  After applying
\(H^k\hookrightarrow H^m\), uniqueness of continuous representatives makes
them equal on \([0,T_*]\).  Their endpoint values are compatible and define
\eqref{u:eq:ier-velocity-endpoint}--%
\eqref{u:eq:ier-mixed-velocity-endpoint}.  Spatial differentiation is
continuous between the corresponding finite Sobolev spaces, so
\(J_{n,\alpha}^*=\partial_x^\alpha V_n^*\).

Corollary~\ref{u:cor:ier-whole-terminal-mixed-jet} gives
\(V_n,V_{n+1}\in L^1(0,T_*;H^m)\) and
\(\partial_tV_n=V_{n+1}\) in distributions.  The Banach-valued fundamental
identity gives
\begin{equation}
 V_n(t)-V_n(s)=\int_s^tV_{n+1}(\tau)\dd\tau
 \quad\text{in }H^m,
 \qquad0\le s\le t\le T_*.
 \label{u:eq:ier-banach-fundamental-identity}
\end{equation}
The continuous representative of \(V_{n+1}\) is bounded by
\(\mathfrak E_{n+1,m}\).  Set \(t=T_*\) in
\eqref{u:eq:ier-banach-fundamental-identity} to prove
\eqref{u:eq:ier-endpoint-velocity-modulus}; spatial differentiation gives
\eqref{u:eq:ier-endpoint-mixed-modulus}.  The remaining size estimates follow
from \eqref{u:eq:ier-row-trace}--\eqref{u:eq:ier-mixed-row-trace}.

The divergence map \(H^{m+1}\to H^m\) is continuous, so
\(\nabla\cdot u_*=0\).  For each \(a,m\), the curves
\(V_0,\ldots,V_{a+1}\) are continuous in \(H^m\), and
\eqref{u:eq:ier-banach-fundamental-identity} identifies the derivative of
\(V_j\) with \(V_{j+1}\).  Induction in \(j\) proves
\eqref{u:eq:ier-complete-left-velocity-jet}.
\end{proof}

\begin{theorem}[Pressure-complete endpoint jet]
\label{u:thm:ier-pressure-complete-endpoint}
\begin{samepage}
Assume the hypotheses of
Theorem~\ref{u:thm:compatible-inverse-intersection}.
For \(n\in\Nzero\), define
\begin{equation}
 Q_n^*:=R_iR_j\sum_{a=0}^{n}\binom na
 V_{a,i}^*V_{n-a,j}^*.
 \label{u:eq:ier-terminal-pressure-row}
\end{equation}
Then \(Q_n^*\in H^\infty(\R^3)\), and, for every \(s\in\Nzero\),
\begin{align}
 \norm{Q_n^*}_{H^s}
 &\le9C_s^{\rm alg}2^n\mathfrak E_{n,s+1}^2,
 \label{u:eq:ier-terminal-pressure-size}\\
 \norm{Q_n(t)-Q_n^*}_{H^s}
 &\le9C_s^{\rm alg}2^{n+1}(T_*-t)\mathfrak E_{n+1,s+1}^2.
 \label{u:eq:ier-terminal-pressure-modulus}
\end{align}
\end{samepage}
For \(\alpha\in\Nzero^3\),
\begin{equation}
 \Pi_{n,\alpha}^*:=\partial_x^\alpha Q_n^*
 =\lim_{t\uparrow T_*}\Pi_{n,\alpha}(t)
 \quad\text{in }H^m(\R^3)\text{ for every }m\in\Nzero.
 \label{u:eq:ier-terminal-mixed-pressure-row}
\end{equation}
More precisely, for every \(m\in\Nzero\),
\begin{align}
 \norm{\Pi_{n,\alpha}^*}_{H^m}
 &\le9C_{m+|\alpha|}^{\rm alg}2^n
 \mathfrak E_{n,m+|\alpha|+1}^2,
 \label{u:eq:ier-terminal-mixed-pressure-size}\\
 \norm{\Pi_{n,\alpha}(t)-\Pi_{n,\alpha}^*}_{H^m}
 &\le9C_{m+|\alpha|}^{\rm alg}2^{n+1}(T_*-t)
 \mathfrak E_{n+1,m+|\alpha|+1}^2.
 \label{u:eq:ier-terminal-mixed-pressure-modulus}
\end{align}
The endpoint rows satisfy, for every \(n,q\in\Nzero\),
\begin{equation}
 V_{n+1}^*
 =\nu\Delta V_n^*
 -\sum_{a=0}^{n}\binom na(V_a^*\cdot\nabla)V_{n-a}^*
 -\nabla Q_n^*
 \quad\text{in }H^q(\R^3).
 \label{u:eq:ier-terminal-velocity-row}
\end{equation}
Among families rooted at \(u_*\) that satisfy
\eqref{u:eq:ier-terminal-pressure-row} and
\eqref{u:eq:ier-terminal-velocity-row} in every finite Sobolev space,
\((V_n^*,Q_n^*)_{n\in\Nzero}\) is the unique pressure-normalized endpoint
jet.  Moreover,
\begin{equation}
 p\in\bigcap_{a,m\in\Nzero}C^a([0,T_*];H^m(\R^3)),
 \qquad
 \partial_t^ap(T_*)=Q_a^*.
 \label{u:eq:ier-complete-left-pressure-jet}
\end{equation}
\end{theorem}

\begin{proof}
Lemmas~\ref{u:lem:sobolev-product} and \ref{u:lem:riesz-leray}, followed by
\eqref{u:eq:ier-row-trace}, give
\begin{align*}
 \norm{Q_n^*}_{H^s}
 &\le9C_s^{\rm alg}\sum_{a=0}^{n}\binom na
 \norm{V_a^*}_{H^{s+1}}\norm{V_{n-a}^*}_{H^{s+1}}\\
 &\le9C_s^{\rm alg}2^n\mathfrak E_{n,s+1}^2.
\end{align*}
For each product in \(Q_n(t)-Q_n^*\), insert
\(V_a^*V_{n-a}(t)\), use
\eqref{u:eq:ier-endpoint-velocity-modulus}, and bound every retained factor by
\(\mathfrak E_{n+1,s+1}\).  The two product differences and
\(\sum_{a=0}^{n}\binom na=2^n\) prove
\eqref{u:eq:ier-terminal-pressure-modulus}.  The Fourier multiplier inequality
\[
 \norm{\partial_x^\alpha f}_{H^m}
 \le\norm f_{H^{m+|\alpha|}}
\]
applied to \eqref{u:eq:ier-terminal-pressure-size} and
\eqref{u:eq:ier-terminal-pressure-modulus} proves
\eqref{u:eq:ier-terminal-mixed-pressure-size} and
\eqref{u:eq:ier-terminal-mixed-pressure-modulus}.  The latter tends to zero
as \(t\uparrow T_*\), which proves
\eqref{u:eq:ier-terminal-mixed-pressure-row} in every \(H^m\).

Fix \(n,q\).  The trace of \(V_n\) in \(H^{q+2}\) carries the diffusion
term in Proposition~\ref{u:prop:formal-mixed-recurrences} to \(H^q\) at
\(T_*\).  The product estimate at order \(q+1\), using velocity convergence
in \(H^{q+2}\), carries every transport term to \(H^q\).  Equation
\eqref{u:eq:ier-terminal-pressure-modulus} at order \(q+1\) carries the pressure
gradient to \(H^q\).  The left side converges to \(V_{n+1}^*\) in \(H^q\),
which proves \eqref{u:eq:ier-terminal-velocity-row}.

The value \(V_0^*=u_*\) fixes \(Q_0^*\), and the velocity recurrence fixes
\(V_1^*\).  Inductively, \(V_0^*,\ldots,V_n^*\) fix \(Q_n^*\), which fixes
\(V_{n+1}^*\).  This proves uniqueness within the displayed recurrence
class.  Extend each \(Q_n\) to \(T_*\) by
\(Q_n(T_*):=Q_n^*\).  On \((0,T_*)\),
\(Q_n=\partial_t^np\), so \(\partial_tQ_n=Q_{n+1}\) in distributions.
Both rows are continuous in every \(H^m\) by
\eqref{u:eq:ier-terminal-pressure-modulus}.  Hence, for
\(0\le s\le t\le T_*\),
\begin{equation}
 Q_n(t)-Q_n(s)=\int_s^tQ_{n+1}(\tau)\dd\tau
 \quad\text{in }H^m(\R^3).
 \label{u:eq:ier-pressure-banach-identity}
\end{equation}
Indeed, the identity holds for \(t<T_*\), and the endpoint bounds permit
passage to \(t=T_*\).  Thus
\(Q_n\in C^1([0,T_*];H^m)\) with derivative \(Q_{n+1}\).  Induction gives
\eqref{u:eq:ier-complete-left-pressure-jet}.
\end{proof}

\begin{lemma}[Late-slice overlap uniqueness]
\label{u:lem:ier-late-slice-overlap}
Assume the hypotheses of
Theorem~\ref{u:thm:compatible-inverse-intersection}.
Define
\begin{equation}
 \delta_0
 :=\min\{1,c_\nu(1+\mathfrak E_{0,3})^{-2}\}>0,
 \label{u:eq:ier-uniform-restart-time}
\end{equation}
where \(c_\nu\) is the constant in Theorem~\ref{u:thm:local-restart}.
For every
\(t_0\in(\max\{0,T_*-\delta_0/2\},T_*)\), let
\[
 (w^{t_0},q^{t_0})
 \in\mathscr M_\nu^3(t_0,t_0+\delta_0;u(t_0)),
 \qquad
 (w^*,q^*)
 \in\mathscr M_\nu^3(T_*,T_*+\delta_0;u_*)
\]
be the unique pairs supplied by Theorem~\ref{u:thm:local-restart}.  Then
\begin{align}
 (w^{t_0},q^{t_0})&=(u,p)
 &&\text{on }[t_0,T_*),
 \label{u:eq:ier-left-overlap}\\
 (w^{t_0},q^{t_0})&=(w^*,q^*)
 &&\text{on }[T_*,t_0+\delta_0].
 \label{u:eq:ier-right-overlap}
\end{align}
\end{lemma}

\begin{proof}
Equation \eqref{u:eq:ier-row-trace} gives
\[
 \norm{u(t_0)}_{H^3}\le\mathfrak E_{0,3},
 \qquad
 \norm{u_*}_{H^3}\le\mathfrak E_{0,3}.
\]
Hence
\[
 \delta_0\le\delta_\nu(u(t_0)),
 \qquad
 \delta_0\le\delta_\nu(u_*),
\]
with \(\delta_\nu\) as in Theorem~\ref{u:thm:local-restart}.
Theorem~\ref{u:thm:local-restart} gives both pairs on the stated intervals.
For each \(t_0<b<T_*\), Lemma~\ref{u:lem:shifted-mild-restriction} places
\((u,p)|_{[t_0,b]}\) and
\((w^{t_0},q^{t_0})|_{[t_0,b]}\) in
\(\mathscr M_\nu^3(t_0,b;u(t_0))\).  Uniqueness in that class gives
\eqref{u:eq:ier-left-overlap}.  Continuity in \(H^3\) gives
\[
 w^{t_0}(T_*)=\lim_{t\uparrow T_*}u(t)=u_*.
\]
Lemma~\ref{u:lem:shifted-mild-restriction} places the restrictions of
\((w^{t_0},q^{t_0})\) and \((w^*,q^*)\) to
\([T_*,t_0+\delta_0]\) in
\(\mathscr M_\nu^3(T_*,t_0+\delta_0;u_*)\).  Uniqueness in that class gives
\eqref{u:eq:ier-right-overlap}.  The right overlap has positive length because
\(t_0+\delta_0>T_*+\delta_0/2\).
\end{proof}

\begin{theorem}[Same-history endpoint restart]
\label{u:thm:ier-same-history-restart}
Assume the hypotheses of
Theorem~\ref{u:thm:compatible-inverse-intersection}.
There is a pressure-normalized classical solution
\((\widetilde u,\widetilde p)\) on
\([0,T_*+\delta_0/2]\times\R^3\) such that
\begin{align}
 (\widetilde u,\widetilde p)|_{[0,T_*)\times\R^3}
 &=(u,p),
 \label{u:eq:ier-extension-identity}\\
 \widetilde u,\widetilde p
 &\in C^\infty([0,T_*+\delta_0/2]\times\R^3),
 \label{u:eq:ier-extension-smoothness}\\
 \partial_t^n\widetilde u(T_*)&=V_n^*,
 \qquad
 \partial_t^n\widetilde p(T_*)=Q_n^*
 \quad(n\in\Nzero).
 \label{u:eq:ier-extension-jet}
\end{align}
If \(0<\eta\le\delta_0/2\) and
\((\widehat u,\widehat p)\) is another pressure-normalized classical
extension satisfying
\begin{equation}
 (\widehat u,\widehat p)|_{[0,T_*)}=(u,p),
 \qquad
 (\widehat u,\widehat p)|_{[T_*,T_*+\eta]}
 \in\mathscr M_\nu^3(T_*,T_*+\eta;u_*),
 \label{u:eq:ier-same-history-uniqueness-class}
\end{equation}
then
\begin{equation}
 (\widehat u,\widehat p)=(\widetilde u,\widetilde p)
 \quad\text{on }[0,T_*+\eta].
 \label{u:eq:ier-same-history-uniqueness}
\end{equation}
\end{theorem}

\begin{proof}
Let \((w^*,q^*)\) be the endpoint solution from
Lemma~\ref{u:lem:ier-late-slice-overlap}, and set
\[
 W_n^+:=\partial_t^nw^*(T_*),
 \qquad
 P_n^+:=\partial_t^nq^*(T_*).
\]
\begin{samepage}
Theorem~\ref{u:thm:local-restart} makes \((w^*,q^*)\) smooth and
pressure-normalized.  Differentiating its equations exactly as in the proof of
Proposition~\ref{u:prop:formal-mixed-recurrences} gives
\begin{align*}
 P_n^+&=R_iR_j\sum_{a=0}^{n}\binom naW_{a,i}^+W_{n-a,j}^+,\\
 W_{n+1}^+&=\nu\Delta W_n^+
 -\sum_{a=0}^{n}\binom na(W_a^+\cdot\nabla)W_{n-a}^+
 -\nabla P_n^+.
\end{align*}
Since \(W_0^+=u_*=V_0^*\), induction with
\eqref{u:eq:ier-terminal-pressure-row} and
\eqref{u:eq:ier-terminal-velocity-row} gives
\begin{equation}
 W_n^+=V_n^*,
 \qquad P_n^+=Q_n^*,
 \qquad n\in\Nzero.
 \label{u:eq:ier-left-right-jet-match}
\end{equation}
\par
\end{samepage}

Define
\[
 (\widetilde u,\widetilde p)(t)
 :=\begin{cases}
 (u,p)(t),&0\le t<T_*,\\
 (w^*,q^*)(t),&T_*\le t\le T_*+\delta_0/2.
 \end{cases}
\]
Equations \eqref{u:eq:ier-complete-left-velocity-jet},
\eqref{u:eq:ier-complete-left-pressure-jet}, and
\eqref{u:eq:ier-left-right-jet-match} match every one-sided time derivative in
every finite Sobolev space.  Sobolev embedding gives
\eqref{u:eq:ier-extension-smoothness}, and the pressure-complete equation holds
at \(T_*\) by \eqref{u:eq:ier-terminal-velocity-row}.  Lemma
\ref{u:lem:ier-late-slice-overlap} identifies the joined pair near \(T_*\) with
the single local history \((w^{t_0},q^{t_0})\), so the gluing is one local
classical history across the endpoint; the piecewise definition gives
\eqref{u:eq:ier-extension-identity}.  Any pair in
\eqref{u:eq:ier-same-history-uniqueness-class} has that datum \(u_*\) on
the right of \(T_*\) and belongs to the shifted mild class containing
\((w^*,q^*)\).  The at-most-one assertion in
Theorem~\ref{u:thm:local-restart} proves
\eqref{u:eq:ier-same-history-uniqueness}.
\end{proof}

The finite rectangle bounds give compatible velocity traces through
\(T_*\) in Theorem~\ref{u:thm:ier-compatible-endpoint-traces}.  The pressure
recurrence fixes their normalized endpoint jet in
Theorem~\ref{u:thm:ier-pressure-complete-endpoint}.  With those endpoint
values, Theorem~\ref{u:thm:ier-same-history-restart} extends the original
history beyond \(T_*\).  A proposed finite maximal time would admit a
classical continuation.  Corollary~\ref{u:cor:ier-maximal-time-contradiction}
applies this contradiction, and Theorem~\ref{u:thm:ier-global-smoothness}
collects the global solution, pressure normalization, and energy identity.

\begin{corollary}[Maximal-time contradiction]
\label{u:cor:ier-maximal-time-contradiction}
The maximal pressure-complete classical history generated by
\((u_0,\nu)\) cannot have \(T_*<\infty\).
\end{corollary}

\begin{proof}
If \(T_*<\infty\), the mixed-index and critical-height descriptions of the
finite rectangles in Theorem~\ref{u:thm:datum-finite-blocks} recover the same
all-orders intersection.  Those rectangles give the trace norms, and the chain from
Theorem~\ref{u:thm:compatible-inverse-intersection} through
Theorem~\ref{u:thm:ier-same-history-restart} produces a pressure-normalized classical
extension to \([0,T_*+\delta_0/2]\), where \(\delta_0>0\).  This contradicts
maximality in Definition~\ref{u:def:formal-maximal-history}.
\end{proof}

\begin{samepage}
\begin{theorem}[Global smoothness on \(\R^3\)]
\label{u:thm:ier-global-smoothness}
For every \(\nu>0\) and every \(u_0\in\HsSigma(\R^3)\), there exists a
pressure-normalized pair
\begin{equation}
 (u,p)\in C^\infty([0,\infty)\times\R^3;\R^3\times\R)
 \label{u:eq:ier-global-class}
\end{equation}
satisfying
\begin{align}
 \partial_tu+(u\cdot\nabla)u+\nabla p&=\nu\Delta u,
 \label{u:eq:ier-global-momentum}\\
 \nabla\cdot u&=0,
 \qquad u(0)=u_0,
 \label{u:eq:ier-global-divergence-datum}\\
 p(t)&=R_iR_j(u_i(t)u_j(t)),
 \qquad \lim_{|x|\to\infty}p(t,x)=0,
 \label{u:eq:ier-global-pressure}\\
 \frac12\norm{u(t)}_{L^2}^2
 +\nu\int_0^t\norm{\nabla u(s)}_{L^2}^2\dd s
 &=\frac12\norm{u_0}_{L^2}^2,
 \qquad t\ge0.
 \label{u:eq:ier-global-energy}
\end{align}
In particular,
\begin{equation}
 \sup_{t\ge0}\norm{u(t)}_{L^2(\R^3)}^2
 \le\norm{u_0}_{L^2(\R^3)}^2<\infty.
 \label{u:eq:ier-global-energy-bound}
\end{equation}
For every finite \(T>0\), its restriction is the unique member of
\(\mathscr M_\nu^3(0,T;u_0)\).
\end{theorem}
\end{samepage}

\begin{proof}
 Let
 \[
  \mathscr T:=\left\{\tau>0:
  \begin{array}{c}
  \text{there exists a pressure-complete classical history}\\
  (\tau,u^\tau,p^\tau)\text{ generated by }(u_0,\nu)
  \end{array}
  \right\},
 \]
 and put
\[
 T_*:=\sup\mathscr T.
\]
 Theorem~\ref{u:thm:local-restart}, applied to \(u_0\in\HsSigma\), supplies a
 shifted mild solution whose persistence statement gives every finite
 velocity and pressure row.  Sobolev embedding and the displayed recurrences
 make this pair a pressure-complete classical history, so \(\mathscr T\) is
 nonempty.
 Lemma~\ref{u:lem:shifted-mild-restriction} admits each strict classical
restriction to the shifted mild class, whose at-most-one property identifies
any two such solutions on their common interval; their union defines a unique
 solution on \([0,T_*)\).  Since \(u_0\in\HsSigma\), the persistence conclusion
of Theorem~\ref{u:thm:local-restart}, applied on overlapping local intervals,
places every strict restriction in the class of
Definition~\ref{u:def:formal-maximal-history}.  An extension beyond \(T_*\)
would restrict to an interval \([0,\tau]\) with \(\tau>T_*\), contrary to
the definition of the supremum.  Thus \((T_*,u,p)\) is maximal.  Corollary
\ref{u:cor:ier-maximal-time-contradiction} gives \(T_*=\infty\).

Fix \(T<\infty\).  Persistence in Theorem~\ref{u:thm:local-restart}, iterated
along the unique history, retains every finite spatial Sobolev order on
\([0,T]\).  Proposition~\ref{u:prop:formal-mixed-recurrences} retains every
finite temporal row.  Lemma~\ref{u:lem:sobolev-embedding} gives
\eqref{u:eq:ier-global-class}; arbitrariness of \(T\) gives the global class.
The pressure belongs to \(H^s(\R^3)\) for every \(s\), and the Fourier
inversion proof of Lemma~\ref{u:lem:sobolev-embedding} places its continuous
representative in \(C_0(\R^3)\), which proves the limit in
\eqref{u:eq:ier-global-pressure}.  Lemma
\ref{u:lem:shifted-mild-restriction} places \((u,p)|_{[0,T]}\) in
\(\mathscr M_\nu^3(0,T;u_0)\).  If
\((\widetilde u,\widetilde p)\in\mathscr M_\nu^3(0,T;u_0)\), the at-most-one
assertion in Theorem~\ref{u:thm:local-restart} gives
\((\widetilde u,\widetilde p)=(u,p)\) on \([0,T]\).  This proves the stated
uniqueness.

Proposition~\ref{u:prop:kinetic-energy} holds on every finite interval.  Since
\(T\) was arbitrary, it gives \eqref{u:eq:ier-global-energy} for every
\(t\ge0\); \eqref{u:eq:ier-global-energy-bound} follows by discarding the
nonnegative dissipation term.
\end{proof}
